\section*{Instructions}
\subsection*{Completing the Assignment}
\begin{itemize}
    \item We provide you the source .tex files that contain the questions. We recommend completing the non-programming part of the assignment in a \LaTeX~editor. The ability to use \LaTeX~is a valuable skill in academia and in industry, and it's generally a clean way to render complicated math. Make sure that your work is formatted correctly. For example, you must submit $\sum_{i=0} x_i$ instead of \text{sum\_\{i=0\} x\_i}.
    \item You may use any \LaTeX~rendering tool you like! \link{https://www.overleaf.com/}{Overleaf} is a great tool, with very few limitations. However, if you want finer control over packages and faster rendering, you may want to download a TeX distribution, a perl binary (if you're on Windows), and the ``Latex Workshop'' extension in VSCode to get live rendering in VSCode.
    \item If you're uncomfortable with using \LaTeX, no worries! You can simply use the pdf version of the non-programming assignment and build your answers in any word or markdown processor you prefer. However, there are some caveats:\begin{enumerate}[label = \arabic*.]
        \item \textbf{We will \underline{NOT} grade handwritten work.} You must use some digital processor to render your work into a cleanly formatted pdf.
        \item It is \textbf{your responsibility} to ensure the questions are correctly numbered. While we will attempt to grade your work accurately, misnumbering may cause misinterpretation.
        \item It is \textbf{your responsibility} to ensure you're answering the question as it is written. We find that students who work without the questions embedded in their document tend to misinterpret the question more often. If you work in a separate document, read carefully and consider copying over the source question, even if its formatting may be lost.
    \end{enumerate}
\end{itemize}
\subsection*{Showing your Work}
\begin{itemize}
    \item When a question asks you to \emph{show your work}, you must provide a justification for your solution. This may include calculations or text-based explanation. Though basic arithmetic may be combined, your work should be thorough and clear enough for a reader to follow your line of reasoning.
    \item When a question asks you to \emph{prove} or \emph{provide a proof of} some statement, you must provide a formal proof. These questions have a higher burden of work, and must constitute a valid inference. Again, while basic arithmetic can be combined, your work should be defensible as a formal proof. Any gaps in logic or unstated assumptions may result in deductions. Be careful using identities you may find online, since they may assume the fact you are trying to prove.
    \item If not otherwise indicated, no justification is needed, just the final answer.
    \item You will never be penalized for showing too much work, but you may be penalized for not showing enough.
    \item If you use an online resource, please cite it. We generally permit most resources, and so, citing it will only be to your benefit in the case of a plagiarism investigation (in the common case where multiple students utilize the same source).
\end{itemize}
\subsection*{Submitting your Work}
\begin{itemize}
    \item We will be using Gradescope for submission and grading of assignments. Complete the assignment, convert your assignment to a pdf, and submit it to the appropriate assignment denoted ``Non-Programming'' on Gradescope.
    \item When submitting your assignment on Gradescope, you are required to map pages of your pdf to subquestions on the homework. You should select a subquestion, then map it to \textbf{every} page which contains any portion of your solution to that subquestion.
    \item \textbf{You must map every page on which you have content that you would like be graded.} During the grading of each particular subquestion, we only see the pages you assign to that subquestion, so you must assign everything you would like us to grade.
    \item \textbf{Improperly mapped questions will receive penalties, up to and including 0 credit for the question.} You are permitted to submit a regrade request in this event; however, the review of such a request is at the sole discretion of the instructional staff and is not likely to be accepted.
    \item Some tips to correctly assign pages:
    \begin{enumerate}[label = \arabic*.]
        \item Underassigning pages is a common mistake. For example, if your answer to question 4.2 starts halfway down on page 16, goes through page 17, and ends with a single line bleeding into page 18, assign pages 16, 17, \textit{and} 18 to question 4.2. Assigning just page 16 may incur a page assignment penalty, or the grader may just interpret your work as incomplete.
        \item Keep in mind that overassigning is just as slowing to the grading process as underassigning is. Please try to keep your assignments tight around the content for the specific subquestion. E.g., assigning every page is just as useless as assigning no pages.
    \end{enumerate}
    \textbf{Again, assign every page that has content relevant to your answer that you would like for us to see while grading. If you don't assign it, we can't see it.}
\end{itemize}